\begin{proof}
\begin{enumerate}[label=\arabic*.,leftmargin=*]
\item Take the canonical target data $(r'_1,r'_2,k'_1,k'_2,l'_0)=(1,-1,0,0,0)$, for which $p'=-1$. Choose $a$ so that $p'=ap$, and an integer $n\equiv k_2\pmod2$. \ref{fam:src:B.2.1} supplies the isomorphism to $X_{c+n}$. The uniqueness is \ref{fam:src:A.5.1}. At fixed $u=e^{-2\pi ic}$, its logarithms differ by integers; the invariant $e^{-\pi i(c+k_2)}$ takes exactly the two square roots of $u$. Both occur, already as $X_c,X_{c+1}$, and \ref{fam:src:A.2.1} distinguishes them.

\item Put $P_{\rm E}=P_{\rm mark}$, the integral isomorphism of
\cref{inv:B1.0}, and retain Engel's $Cf_4=e_1+e_4$.
Besides the intertwining identities in that proposition, multiplication gives
$P_{\rm E}d=\delta_0$, $\psi P_{\rm E}=-\gamma$,
$P_{\rm E}\varepsilon=-e_4$ and
$P_{\rm E}\varepsilon'=\delta_0-Cf_4$.
They also agree with the complex periods: at $z_1$ the complex-linear map
$(\zeta_1,\zeta_2)\mapsto(-\zeta_1/(\rho+1),\zeta_2+\zeta_1/3)$ carries $\Pi_cP_{\rm E}^{-1}$ to Engel's columns $(1,0),(\rho-1,0),((1+\rho)/3,1/3),(0,-\kappa_1)$. At $z_2$ his $Cf_4$ has preimage $d-\varepsilon'$ and fixed period $1-\kappa_2$. Thus the comparison fixes the positive circle class and the local complex structures, in addition to their monodromy matrices.

\item The inverse-generator convention of \cref{inv:B1.0} gives $a_{j-1}=-P_{\rm E}v_j/m_j$ in Engel's lattice marking (and its period image in the torus). Integer shift gives $\Pi_{c+1}v_j=\Pi_c(v_j+\gamma(v_j)d)$, with $\gamma(v_1)=1,\gamma(v_2)=-1$, proving the two signs. A positive $u$ circuit changes $\log u$ by $2\pi i$ and hence changes $c$ by $-1$, exchanging the square roots. This gives the claimed permutation with its orientation.
\end{enumerate}
\end{proof}